\documentclass[10pt,a4paper]{amsart}
\usepackage[margin=19mm]{geometry}
\usepackage{amsmath,amssymb,mathtools}
\usepackage[scr=rsfs,cal=euler]{mathalfa}
\usepackage{xfrac}
\usepackage[unicode,hidelinks,bookmarks=false]{hyperref}
\newcommand{\ie}{i.e.\ }
\newcommand{\cf}{cf.\ }
\newcommand{\resp}{respectively\ }
\newcommand{\Subsection}{\subsection}
\providecommand{\nobreakdash}{\nobreak\hbox{-}}
\setlength{\emergencystretch}{2em}
\multlinegap0pt
\usepackage{amssymb}
\usepackage{mathtools}
\usepackage{xcolor}
\newtheorem{lemma}{Lemma}
\newtheorem{theorem}{Theorem}
\newcommand{\Edist}[2]      {\mm{\|{#1}-{#2}\|}}
\newcommand{\Overlay}[3]    {\mm{{\rm Vor}_{#1}{({#2}\mid{#3})}}}
\newcommand{\Rspace}        {\mm{{\mathbb R}}}
\newcommand{\Union}[2]      {\mm{\rm Union}{({#1},{#2})}}
\newcommand{\Voronoi}[2]    {\mm{{\rm Vor}_{#1}{({#2})}}}
\newcommand{\card}[1]       {\mm{{\#}{#1}}}
\newcommand{\length}[1]     {\mm{\rm length}{[{#1}]}}
\newcommand {\mm}[1]   {\ifmmode{#1}\else{\mbox{\(#1\)}}\fi}
\title{On the Size of Chromatic Delaunay Mosaics}
\author{Ranita Biswas, Sebastiano Cultrera di Montesano, Ond\v{r}ej Draganov, Herbert Edelsbrunner, Morteza Saghafian}
\date{Source: arXiv:2212.03121v1 (CC BY 4.0)}
\begin{document}
\maketitle

\noindent\emph{Source and licence.} On the Size of Chromatic Delaunay Mosaics. Version arXiv:2212.03121v1 (CC BY 4.0), distributed by arXiv under the Creative Commons Attribution 4.0 International licence, http://creativecommons.org/licenses/by/4.0/, as recorded in the arXiv licence metadata for this exact version and shown on the abstract page https://arxiv.org/abs/2212.03121v1. Full licence notice: \texttt{licenses/arxiv-2212.03121-CC-BY-4.0.txt}.\par\smallskip

\noindent\textit{Editorial scope.} Counted unit: the article's theorem thm:overlay\_size\_for\_dense\_set (source0.tex lines 645--648). The article's complete original proof of it is delivered with it (source0.tex lines 649--705, one \texttt{proof} environment of 57 lines). No mathematics is written, repaired or re-derived: the statement and the proof are the article's own lines, in the article's own notation, and no retained line is substituted or re-flowed.  Retained uncounted as the source-local context the proof needs: lines 567--580; lines 616--621.  Nothing was omitted from any retained span, so the package carries no omission record.  No retained line contains a citation, so the wrapper carries no bibliography and no bibliography entry.  No retained line contains a figure environment, an image-inclusion command or drawing code, so no image is delivered and the package declares no asset dependency.  Editorial formatting only, in addition: an AMS \texttt{amsart} preamble that restates the article's own declarations and macros used by the retained lines, namely the environments \texttt{lemma}, \texttt{theorem} on separate counters, together with the standard packages those lines need.  The retained environments print in this document as Lemma 1, Theorem 1 in order of appearance, whereas the article prints the same items with the numbering of its own full document.  The complete native source of the article as distributed in the arXiv e-print for this exact version is bundled unchanged as \texttt{sources/134549/source0.tex}.
\begin{lemma}[Boundary of Union of Disks]
  \label{lem:boundary_of_union_of_disks}
  Let $S$ be a line segment of length $L$, $Q$ a square with sides of length $L$, and $\bar{Q}$ the closed complement of $Q$.
  For $Y \in \{S, Q, \bar{Q} \}$, let $\varrho_Y \colon Y \to \Rspace$ be non-negative and $R_Y = \max_{x \in Y} \varrho_Y (x)$.
  Then 
  \begin{align}
    \length{\partial \Union{S}{\varrho_S}} &< 4L+8R_S, 
      \label{eqn:lengthS} \\ 
    \length{\partial \Union{Q}{\varrho_Q}} &< 8L+8R_Q,
      \label{eqn:lengthQ} \\ 
    \length{\partial \Union{\bar{Q}}{\varrho_{\bar{Q}}}} &< 8L.
      \label{eqn:lengthbarQ}
  \end{align}
\end{lemma}

\begin{lemma}[Counting Points and Crossings]
  \label{lem:counting_points_and_crossings}
  Let $S$ be a line segment, $Q$ a square, $\bar{Q}$ the closed complement of $Q$, and $B \subseteq \Rspace^2$ finite.
  For $Y \in \{ S, Q, \bar{Q} \}$, let $\varrho_Y \colon Y \to \Rspace$ be defined by $\varrho_Y (x) = \min_{a \in B} \Edist{x}{a}$, and write $B_Y = B \cap \partial \Union{Y}{\varrho_Y}$.
  Then the number of edges of $\Voronoi{}{B}$ that have a non-empty intersection with $S$, $\partial Q$, $\partial \bar{Q}$ is bounded from above by $\card{B_S}$, $\card{B_Q}$, $\card{B_{\bar{Q}}}$, respectively.
\end{lemma}

\begin{theorem}[Overlay size for Dense set]\label{thm:overlay_size_for_dense_set}
  Let $A \subseteq \Rspace^2$ be finite with density $m$, write $n = \card{A}$, let $\sigma = \{0,1,\ldots,s\}$, and write $A_j = \chi^{-1} (j)$, in which $\chi \colon A \to \sigma$ is any coloring.
  Then the number of regions in $\Overlay{}{A_j}{j \in \sigma}$ is $O(s^2 m^2)$.
\end{theorem}

\begin{proof}
  We will show that the number of crossings between the edges of any two mono-chromatic Voronoi tessellations is $O(m^2)$.
  There are $\binom{s+1}{2} \leq s^2$ pairs of colors and therefore $O(s^2 m^2)$ crossings in total.
  The number of regions in the overlay is the number of regions in the $s+1$ mono-chromatic Voronoi tessellations, which is $n = \card{A}$, plus twice the number of crossings.
  Since $n = O(m^2)$, this implies that the number of regions is $O(s^2 m^2)$.
  
  \medskip
  For the remainder of this proof, we fix two colors, $0$ and $1$, we assume that the minimum distance between points in $A$ is $1$, so the maximum distance is $m$.
  Observe that there is a square with sides of length $m$ that contains $A$, and therefore $A_0$ and $A_1$.
  If there is at least one point each of $A_0$ and $A_1$ in the square, then we subdivide it into four equally large squares.
  We recursively subdivide each of these squares independently until we arrive at squares that contain points of at most one of these two colors.
  By choosing the initial square judiciously, we may assume that no point of $A$ lies on the boundary of any of these squares.
  Since subdivision does not alter the total area, the sum of areas of these squares is $m^2$.
  
  Let $Q$ be a square in this subdivision, write $L$ for the length of its sides, and assume that it contains no points of $A_0$.
  Let $Q'$ be the parent square of four times the area, which, by construction, contains at least one point of $A_0$ and at least one point of $A_1$.
  Let $\varrho_Q \colon Q \to \Rspace$ be defined by $\varrho_Q (x) = \min_{a \in A_0} \Edist{x}{a}$.
  Since $Q'$ contains at least one point of $A_0$, we have $R_Q = \max_{x \in Q} \varrho_Q (x) < 2 \sqrt{2} L$.
  Recall that $\Union{Q}{\varrho_Q}$ is the union of closed disks with centers $x$ and radii $\varrho_Q (x)$.
  By Lemma~\ref{lem:boundary_of_union_of_disks}, the length of the boundary satisfies
  \begin{align}
    \length{\partial \Union{Q}{\varrho_Q}}  &<  8L + 8R_Q  < (8 + 16 \sqrt{2}) L  <  31 L .
  \end{align}
  Since any two points of $A_0$ are at least a distance $1$ apart, this implies that there are fewer than $31 L$ points of $A_0$ on the boundary of the union of disks.
  By Lemma~\ref{lem:counting_points_and_crossings}, fewer than $31 L$ edges in $\Voronoi{}{A_0}$ cross the sides of $Q$.
  Since no point of $A_0$ is inside $Q$, the edges of $\Voronoi{}{A_0}$ inside $Q$ do not form cycles, so more than half of them cross the sides of $Q$.
  It follows that fewer than $62 L$ edges of $\Voronoi{}{A_0}$ have a non-empty intersection with $Q$.
  %
  Let $S$ be the intersection of one such edge with $Q$, which is either the entire edge or a connected piece of it.
  The length of $S$ is at most $\sqrt{2} L$.
  Let $\varrho_S \colon S \to \Rspace$ be defined by $\varrho_S (x) = \min_{a \in A_1} \Edist{x}{a}$.
  The maximum such distance satisfies $R_S = \max_{x \in S} \varrho_S (x) < 2 \sqrt{2} L$.
  By~Lemma~\ref{lem:boundary_of_union_of_disks}, the length of the boundary satisfies
  \begin{align}
    \length{\partial \Union{S}{\varrho_S}}  &<  4L+8R_S  <  (4 + 16 \sqrt{2}) L  <  27 L .
  \end{align}
  Since any two points of $A_1$ are at least a distance $1$ apart, this implies that fewer than $27 L$ points of $A_1$ lie on the boundary of the union of disks.
  By Lemma~\ref{lem:counting_points_and_crossings}, fewer than $27 L$ edges of $\Voronoi{}{A_1}$ cross $S$.
  %
  Multiplying with the number of edges in $\Voronoi{}{A_0}$ inside $Q$, we get fewer than $62 L \cdot 27 L = 1674 L^2$ crossings.
  This is only a constant times the area of $Q$.
  Taking the sum over all squares in the subdivision, we thus get fewer than $1674 m^2$ crossings between edges of $\Voronoi{}{A_0}$ and $\Voronoi{}{A_1}$ inside the initial square.
  
  \medskip
  It remains to bound the number of crossings outside the initial square.
  Let $\bar{Q}$ be the complement of the initial square, which we recall has sides of length $m$.
  Let $\varrho_j \colon \bar{Q} \to \Rspace$ be defined by $\varrho_j (x) = \min_{a \in A_j} \Edist{x}{a}$, for $j = 0, 1$.
  By Lemma~\ref{lem:boundary_of_union_of_disks}, we have
  \begin{align}
    \length{\partial \Union{\bar{Q}}{\varrho_j}}  &<  8m .
  \end{align}
  By Lemma~\ref{lem:counting_points_and_crossings}, fewer than $8m$ edges of $\Voronoi{}{A_j}$ cross the sides of $\bar{Q}$.
  Since there are no points of $A_j$ in $\bar{Q}$,
  this implies that $\Voronoi{}{A_j}$ has fewer than $16 m$ edges with non-empty intersection with $\bar{Q}$.
  Even if every such edge of $\Voronoi{}{A_0}$ crossed every such edge of $\Voronoi{}{A_1}$, we still have fewer than $16 m \cdot 16 m = 256 m^2$ crossings outside $Q$.
  Adding the numbers of crossings inside and outside $Q$, we get fewer than $1930 m^2$ crossings altogether.
\end{proof}

\end{document}
